\documentclass[11pt,a4paper]{article}

%===========================================================
% PACKAGES
%===========================================================
\usepackage[utf8]{inputenc}
\usepackage[T1]{fontenc}
\usepackage[french]{babel}
\usepackage{lmodern}
\usepackage[a4paper,margin=2.2cm]{geometry}
\usepackage{amsmath,amssymb,amsthm,mathtools}
\usepackage{fancyhdr}
\usepackage{lastpage}
\usepackage{tikz}
\usepackage[most]{tcolorbox}
\usepackage{enumitem}
\usepackage{array,tabularx,booktabs}
\usepackage{hyperref}
\usepackage{xcolor}
\usepackage{microtype}
\usepackage{multicol}

%===========================================================
% MISE EN PAGE
%===========================================================
\setlength{\headheight}{15pt}
\pagestyle{fancy}
\fancyhf{}
\fancyhead[L]{Classe préparatoire scientifique}
\fancyhead[C]{Logique et raisonnement}
\fancyhead[R]{Page \thepage/\pageref{LastPage}}
\fancyfoot[C]{Pierre Garcia-Buscaylet -- Polycopié de cours}

\hypersetup{
  colorlinks=true,
  linkcolor=blue!50!black,
  urlcolor=blue!50!black,
  citecolor=blue!50!black
}

\setlist[itemize]{label=--,leftmargin=1.1cm}
\setlist[enumerate]{leftmargin=1.2cm}
\renewcommand{\arraystretch}{1.25}

%===========================================================
% COULEURS ET BOITES
%===========================================================
\definecolor{bleuprepa}{RGB}{20,55,110}
\definecolor{vertprepa}{RGB}{20,105,70}
\definecolor{rougeprepa}{RGB}{160,35,35}
\definecolor{orangeprepa}{RGB}{180,100,20}
\definecolor{grisclair}{RGB}{245,247,250}

\tcbset{
  enhanced,
  breakable,
  sharp corners=downhill,
  boxrule=0.7pt,
  colback=white,
  fonttitle=\bfseries
}

\newtcbtheorem[number within=section]{definitionbox}{Définition}{colframe=bleuprepa,colback=blue!3!white}{def}
\newtcbtheorem[number within=section]{propositionbox}{Proposition}{colframe=vertprepa,colback=green!3!white}{prop}
\newtcbtheorem[number within=section]{theorembox}{Théorème}{colframe=bleuprepa!80!black,colback=blue!4!white}{thm}
\newtcbtheorem[number within=section]{lemmebox}{Lemme}{colframe=vertprepa!80!black,colback=green!3!white}{lem}
\newtcbtheorem[number within=section]{exemplebox}{Exemple corrigé}{colframe=black!70,colback=grisclair}{ex}
\newtcbtheorem[number within=section]{exercicebox}{Exercice}{colframe=orangeprepa,colback=orange!4!white}{exo}

\newtcolorbox{methodebox}[1][]{title={Méthode},colframe=bleuprepa,colback=blue!2!white,#1}
\newtcolorbox{retenirbox}[1][]{title={À retenir},colframe=vertprepa,colback=green!4!white,#1}
\newtcolorbox{erreurbox}[1][]{title={Erreur classique},colframe=rougeprepa,colback=red!3!white,#1}
\newtcolorbox{preuvebox}[1][]{title={Démonstration},colframe=black!70,colback=white,#1}
\newtcolorbox{remarquebox}[1][]{title={Remarque pédagogique},colframe=black!55,colback=black!2!white,#1}

%===========================================================
% COMMANDES
%===========================================================
\newcommand{\V}{\mathrm{V}}
\newcommand{\F}{\mathrm{F}}
\newcommand{\ssi}{\Longleftrightarrow}
\newcommand{\implique}{\Longrightarrow}
\newcommand{\R}{\mathbb{R}}
\newcommand{\N}{\mathbb{N}}
\newcommand{\Z}{\mathbb{Z}}
\newcommand{\Q}{\mathbb{Q}}
\newcommand{\C}{\mathbb{C}}
\newcommand{\calP}{\mathcal{P}}
\newcommand{\id}{\mathrm{id}}
\newcommand{\Card}{\mathrm{Card}}

\title{\vspace{-1cm}\Huge\bfseries Logique, raisonnement et rédaction mathématique\\[0.2cm]
\Large Cours complet pour classes préparatoires scientifiques}
\author{Support professeur -- Niveau MPSI/PCSI/MP2I/PSI}
\date{\today}

\begin{document}
\maketitle
\thispagestyle{empty}

\begin{center}
\begin{tcolorbox}[width=0.92\textwidth,colframe=bleuprepa,colback=blue!2!white,title=Objectif du polycopié]
Ce cours pose les règles du jeu des mathématiques supérieures : ce qu'est une proposition, comment se manipulent les connecteurs logiques, comment se lisent et se nient les quantificateurs, comment se rédigent les principaux types de démonstration, et comment relier logique, ensembles, applications et rédaction rigoureuse. L'objectif n'est pas seulement de connaître des symboles, mais de savoir produire une preuve claire, complète et vérifiable.
\end{tcolorbox}
\end{center}

\tableofcontents
\newpage

%===========================================================
\section*{Préambule : pourquoi commencer par la logique ?}
\addcontentsline{toc}{section}{Préambule : pourquoi commencer par la logique ?}

En mathématiques, une affirmation n'est pas acceptée parce qu'elle semble vraie sur quelques exemples, parce qu'elle est visible sur une figure, ou parce qu'elle est plausible. Elle est acceptée lorsqu'elle est démontrée à partir de définitions, d'axiomes ou de résultats déjà établis.

La logique fournit le langage minimal permettant de savoir :
\begin{itemize}
  \item ce que signifie exactement une phrase mathématique ;
  \item comment transformer une affirmation sans changer son sens ;
  \item comment nier correctement une assertion ;
  \item comment construire une preuve ;
  \item comment détecter une erreur de raisonnement.
\end{itemize}

\begin{retenirbox}
Faire des mathématiques, ce n'est pas manipuler des symboles au hasard : c'est produire des chaînes d'implications dont chaque maillon est justifié. La logique est donc l'outil qui permet de contrôler la validité de ces chaînes.
\end{retenirbox}

%===========================================================
\section{Propositions, prédicats et vérité}

\subsection{Proposition mathématique}

\begin{definitionbox}{Proposition}{proposition}
Une \emph{proposition} est une phrase mathématique à laquelle on peut attribuer une valeur de vérité : vraie ou fausse.
\end{definitionbox}

\begin{exemplebox}{Exemples et non-exemples}{prop-exemples}
Les phrases suivantes sont des propositions :
\begin{enumerate}
  \item $2+3=5$, qui est vraie ;
  \item $7$ est un nombre pair, qui est fausse ;
  \item pour tout $x\in\R$, $x^2\geq 0$, qui est vraie.
\end{enumerate}
La phrase « $x^2=1$ » n'est pas une proposition tant que $x$ n'est pas fixé ou quantifié. Elle devient une proposition si l'on écrit par exemple :
\[
\forall x\in\R,\ x^2=1,
\]
qui est fausse, ou
\[
\exists x\in\R,\ x^2=1,
\]
qui est vraie.
\end{exemplebox}

\begin{definitionbox}{Prédicat}{predicat}
Un \emph{prédicat} est une phrase dépendant d'une ou plusieurs variables, qui devient une proposition lorsque les variables sont fixées ou quantifiées. On note souvent un prédicat $P(x)$, $P(x,y)$, etc.
\end{definitionbox}

\begin{exemplebox}{Un prédicat}{predicat-exemple}
Soit $P(x)$ la phrase : « $x^2-1=0$ ». Alors :
\begin{itemize}
  \item $P(1)$ est vraie ;
  \item $P(2)$ est fausse ;
  \item $\exists x\in\R, P(x)$ est vraie ;
  \item $\forall x\in\R, P(x)$ est fausse.
\end{itemize}
\end{exemplebox}

\begin{erreurbox}
Ne pas confondre une expression ouverte, comme $x>0$, avec une proposition. Tant que l'on ne sait pas qui est $x$, ou tant que $x$ n'est pas quantifié, la phrase n'a pas de valeur de vérité globale.
\end{erreurbox}

\subsection{Syntaxe et sémantique}

\begin{definitionbox}{Syntaxe et sémantique}{syntaxe-semantique}
La \emph{syntaxe} concerne la forme d'une phrase : symboles utilisés, ordre des quantificateurs, parenthèses, variables libres ou liées. La \emph{sémantique} concerne le sens et la valeur de vérité de la phrase.
\end{definitionbox}

Deux phrases peuvent avoir une syntaxe différente mais la même signification. Par exemple :
\[
\forall x\in\R,\ x^2\geq 0
\qquad\text{et}\qquad
\text{tout réel a un carré positif ou nul.}
\]
À l'inverse, deux phrases très proches syntaxiquement peuvent avoir des sens très différents :
\[
\forall x\in\R,\exists y\in\R,\ y>x
\]
est vraie, tandis que
\[
\exists y\in\R,\forall x\in\R,\ y>x
\]
est fausse.

\subsection{Exercices d'application}

\begin{exercicebox}{Identifier propositions et prédicats}{exo-prop}
Parmi les phrases suivantes, dire lesquelles sont des propositions. Lorsqu'il s'agit d'un prédicat, préciser les variables libres.
\begin{enumerate}
  \item $3^2=9$.
  \item $x+1=0$.
  \item Pour tout $n\in\N$, $n^2+n$ est pair.
  \item Il existe $x\in\R$ tel que $x^2+1=0$.
  \item $x<y$.
\end{enumerate}
\end{exercicebox}

\begin{exercicebox}{Transformer en proposition}{exo-transformer}
Transformer les phrases suivantes en propositions vraies, puis en propositions fausses, en ajoutant des quantificateurs ou des conditions.
\begin{enumerate}
  \item $x^2\geq 0$.
  \item $x^2=4$.
  \item $n$ est divisible par $2$.
\end{enumerate}
\end{exercicebox}

%===========================================================
\section{Connecteurs logiques et tables de vérité}

\subsection{Négation}

\begin{definitionbox}{Négation}{negation}
Si $P$ est une proposition, sa négation est la proposition notée $\neg P$ ou $\overline P$, qui est vraie exactement lorsque $P$ est fausse.
\end{definitionbox}

\begin{center}
\begin{tabular}{c|c}
$P$ & $\neg P$ \\
\hline
$\V$ & $\F$ \\
$\F$ & $\V$
\end{tabular}
\end{center}

\begin{exemplebox}{Nier une phrase simple}{nier-simple}
La négation de « $x\geq 3$ » est « $x<3$ ». La négation de « $x=0$ » est « $x\neq 0$ ». La négation de « $x<0$ » est « $x\geq 0$ ».
\end{exemplebox}

\begin{erreurbox}
La négation de « tous les élèves ont compris » n'est pas « aucun élève n'a compris », mais « au moins un élève n'a pas compris ».
\end{erreurbox}

\subsection{Conjonction : « et »}

\begin{definitionbox}{Conjonction}{conjonction}
Si $P$ et $Q$ sont deux propositions, la conjonction $P\land Q$ signifie « $P$ et $Q$ ». Elle est vraie lorsque $P$ et $Q$ sont toutes les deux vraies.
\end{definitionbox}

\begin{center}
\begin{tabular}{c|c|c}
$P$ & $Q$ & $P\land Q$ \\
\hline
$\V$ & $\V$ & $\V$ \\
$\V$ & $\F$ & $\F$ \\
$\F$ & $\V$ & $\F$ \\
$\F$ & $\F$ & $\F$
\end{tabular}
\end{center}

\subsection{Disjonction inclusive : « ou »}

\begin{definitionbox}{Disjonction inclusive}{disjonction}
Si $P$ et $Q$ sont deux propositions, la disjonction $P\lor Q$ signifie « $P$ ou $Q$ », au sens inclusif : elle est vraie lorsqu'au moins l'une des deux propositions est vraie.
\end{definitionbox}

\begin{center}
\begin{tabular}{c|c|c}
$P$ & $Q$ & $P\lor Q$ \\
\hline
$\V$ & $\V$ & $\V$ \\
$\V$ & $\F$ & $\V$ \\
$\F$ & $\V$ & $\V$ \\
$\F$ & $\F$ & $\F$
\end{tabular}
\end{center}

\begin{remarquebox}
Dans la langue courante, « ou » est parfois exclusif : « fromage ou dessert » signifie souvent « l'un ou l'autre, mais pas les deux ». En mathématiques, sauf précision contraire, le « ou » est inclusif.
\end{remarquebox}

\subsection{Implication}

\begin{definitionbox}{Implication}{implication}
L'implication $P\Rightarrow Q$ se lit « si $P$, alors $Q$ ». Elle est fausse uniquement dans le cas où $P$ est vraie et $Q$ est fausse.
\end{definitionbox}

\begin{center}
\begin{tabular}{c|c|c}
$P$ & $Q$ & $P\Rightarrow Q$ \\
\hline
$\V$ & $\V$ & $\V$ \\
$\V$ & $\F$ & $\F$ \\
$\F$ & $\V$ & $\V$ \\
$\F$ & $\F$ & $\V$
\end{tabular}
\end{center}

\begin{retenirbox}
Une implication $P\Rightarrow Q$ ne dit pas que $P$ est vraie. Elle dit seulement que, si $P$ est vraie, alors $Q$ doit l'être. Lorsque $P$ est fausse, l'implication est vraie par défaut : on parle de vérité vacue.
\end{retenirbox}

\begin{exemplebox}{Implication vraie avec hypothèse fausse}{imp-vacue}
La proposition suivante est vraie :
\[
\forall x\in\R,\quad x^2<0 \Rightarrow x=17.
\]
En effet, il n'existe aucun réel $x$ tel que $x^2<0$. L'hypothèse est donc toujours fausse, et l'implication ne peut jamais être mise en défaut.
\end{exemplebox}

\subsection{Réciproque, contraposée et négation d'une implication}

\begin{definitionbox}{Réciproque et contraposée}{reciproque-contraposee}
Soit une implication $P\Rightarrow Q$.
\begin{itemize}
  \item Sa \emph{réciproque} est $Q\Rightarrow P$.
  \item Sa \emph{contraposée} est $\neg Q\Rightarrow \neg P$.
  \item Sa \emph{négation} est $P\land \neg Q$.
\end{itemize}
\end{definitionbox}

\begin{propositionbox}{Équivalence entre implication et contraposée}{equiv-contraposee}
Les propositions $P\Rightarrow Q$ et $\neg Q\Rightarrow\neg P$ sont logiquement équivalentes.
\end{propositionbox}

\begin{preuvebox}
On compare les tables de vérité :
\[
\begin{array}{c|c|c|c|c|c}
P&Q&P\Rightarrow Q&\neg Q&\neg P& \neg Q\Rightarrow\neg P\\
\hline
\V&\V&\V&\F&\F&\V\\
\V&\F&\F&\V&\F&\F\\
\F&\V&\V&\F&\V&\V\\
\F&\F&\V&\V&\V&\V
\end{array}
\]
Les colonnes de $P\Rightarrow Q$ et $\neg Q\Rightarrow\neg P$ coïncident. Les deux propositions sont donc équivalentes.
\end{preuvebox}

\begin{erreurbox}
La réciproque $Q\Rightarrow P$ n'est pas équivalente à $P\Rightarrow Q$. Par exemple, si $n$ est divisible par $4$, alors $n$ est pair. La réciproque « si $n$ est pair, alors $n$ est divisible par $4$ » est fausse.
\end{erreurbox}

\subsection{Équivalence}

\begin{definitionbox}{Équivalence}{equivalence}
L'équivalence $P\Longleftrightarrow Q$ signifie que $P$ et $Q$ ont la même valeur de vérité. Elle est équivalente à la conjonction de deux implications :
\[
P\Longleftrightarrow Q \quad\equiv\quad (P\Rightarrow Q)\land(Q\Rightarrow P).
\]
\end{definitionbox}

\begin{center}
\begin{tabular}{c|c|c}
$P$ & $Q$ & $P\Longleftrightarrow Q$ \\
\hline
$\V$ & $\V$ & $\V$ \\
$\V$ & $\F$ & $\F$ \\
$\F$ & $\V$ & $\F$ \\
$\F$ & $\F$ & $\V$
\end{tabular}
\end{center}

\begin{methodebox}
Pour démontrer $P\Longleftrightarrow Q$, on démontre souvent séparément :
\begin{enumerate}
  \item $P\Rightarrow Q$ ;
  \item $Q\Rightarrow P$.
\end{enumerate}
On peut aussi utiliser une chaîne d'équivalences, mais seulement si chaque étape est réellement réversible.
\end{methodebox}

\subsection{Lois logiques fondamentales}

\begin{propositionbox}{Lois de De Morgan}{de-morgan}
Pour toutes propositions $P$ et $Q$ :
\[
\neg(P\land Q)\Longleftrightarrow (\neg P)\lor(\neg Q),
\]
\[
\neg(P\lor Q)\Longleftrightarrow (\neg P)\land(\neg Q).
\]
\end{propositionbox}

\begin{preuvebox}
Démontrons la première loi par table de vérité :
\[
\begin{array}{c|c|c|c|c|c|c}
P&Q&P\land Q&\neg(P\land Q)&\neg P&\neg Q&(\neg P)\lor(\neg Q)\\
\hline
\V&\V&\V&\F&\F&\F&\F\\
\V&\F&\F&\V&\F&\V&\V\\
\F&\V&\F&\V&\V&\F&\V\\
\F&\F&\F&\V&\V&\V&\V
\end{array}
\]
Les colonnes de $\neg(P\land Q)$ et $(\neg P)\lor(\neg Q)$ sont identiques. La seconde loi se démontre de même.
\end{preuvebox}

\begin{propositionbox}{Distributivité}{distributivite-logique}
Pour toutes propositions $P,Q,R$ :
\[
P\land(Q\lor R)\Longleftrightarrow (P\land Q)\lor(P\land R),
\]
\[
P\lor(Q\land R)\Longleftrightarrow (P\lor Q)\land(P\lor R).
\]
\end{propositionbox}

\begin{retenirbox}
Les connecteurs logiques obéissent à des règles proches de celles des opérations ensemblistes. Cette analogie sera systématisée avec les ensembles : $\land$ correspond à l'intersection, $\lor$ à la réunion, et la négation au complémentaire.
\end{retenirbox}

\subsection{Exercices d'application et bilan}

\begin{exercicebox}{Tables de vérité}{exo-table-verite}
Construire les tables de vérité des propositions suivantes et dire lesquelles sont des tautologies :
\begin{enumerate}
  \item $P\Rightarrow(P\lor Q)$ ;
  \item $(P\land Q)\Rightarrow P$ ;
  \item $(P\Rightarrow Q)\Longleftrightarrow(\neg P\lor Q)$ ;
  \item $(P\Rightarrow Q)\Longleftrightarrow(Q\Rightarrow P)$.
\end{enumerate}
\end{exercicebox}

\begin{exercicebox}{Nier correctement}{exo-negations}
Nier les propositions suivantes :
\begin{enumerate}
  \item $x>0$ et $x<1$.
  \item $n$ est pair ou divisible par $3$.
  \item Si $x>2$, alors $x^2>4$.
  \item $P\Longleftrightarrow Q$.
\end{enumerate}
\end{exercicebox}

%===========================================================
\section{Quantificateurs}

\subsection{Quantificateur universel}

\begin{definitionbox}{Quantificateur universel}{forall}
La proposition
\[
\forall x\in E,\ P(x)
\]
se lit « pour tout $x$ appartenant à $E$, $P(x)$ est vraie ». Elle signifie que chaque élément de $E$ vérifie la propriété $P$.
\end{definitionbox}

\begin{exemplebox}{Phrase universelle}{ex-forall}
La proposition
\[
\forall x\in\R,\quad x^2\geq 0
\]
est vraie. Pour la démontrer, on prend un réel arbitraire $x$, puis on prouve que $x^2\geq 0$.
\end{exemplebox}

\subsection{Quantificateur existentiel}

\begin{definitionbox}{Quantificateur existentiel}{exists}
La proposition
\[
\exists x\in E,\ P(x)
\]
se lit « il existe au moins un $x$ appartenant à $E$ tel que $P(x)$ soit vraie ».
\end{definitionbox}

\begin{definitionbox}{Existence unique}{existsunique}
La proposition
\[
\exists ! x\in E,\ P(x)
\]
se lit « il existe un unique $x$ appartenant à $E$ tel que $P(x)$ soit vraie ». Elle signifie :
\[
\exists x\in E,\ P(x)
\quad\text{et}\quad
\forall y\in E,\ P(y)\Rightarrow y=x.
\]
\end{definitionbox}

\begin{methodebox}
Pour prouver une existence unique, on sépare généralement deux étapes :
\begin{enumerate}
  \item \textbf{Existence :} on construit un objet qui convient, ou on démontre qu'il existe.
  \item \textbf{Unicité :} on suppose que deux objets conviennent et on montre qu'ils sont égaux.
\end{enumerate}
\end{methodebox}

\subsection{Ordre des quantificateurs}

L'ordre des quantificateurs est essentiel. Les deux propositions
\[
\forall x\in E,\exists y\in F,\ P(x,y)
\]
et
\[
\exists y\in F,\forall x\in E,\ P(x,y)
\]
ne disent pas la même chose.

\begin{exemplebox}{Ordre des quantificateurs}{ordre-quant}
Dans $\R$, la proposition
\[
\forall x\in\R,\exists y\in\R,\ y>x
\]
est vraie : pour chaque $x$, on peut prendre $y=x+1$. Mais la proposition
\[
\exists y\in\R,\forall x\in\R,\ y>x
\]
est fausse : il n'existe pas de réel plus grand que tous les réels.
\end{exemplebox}

\begin{erreurbox}
Intervertir deux quantificateurs de nature différente change généralement le sens de la phrase. En revanche, deux quantificateurs de même nature peuvent être permutés :
\[
\forall x\forall y\ P(x,y) \Longleftrightarrow \forall y\forall x\ P(x,y),
\]
\[
\exists x\exists y\ P(x,y) \Longleftrightarrow \exists y\exists x\ P(x,y).
\]
\end{erreurbox}

\subsection{Négation des propositions quantifiées}

\begin{theorembox}{Négation des quantificateurs}{neg-quant}
Soit $P(x)$ un prédicat sur un ensemble $E$. Alors :
\[
\neg\left(\forall x\in E,\ P(x)\right)
\Longleftrightarrow
\exists x\in E,\ \neg P(x),
\]
\[
\neg\left(\exists x\in E,\ P(x)\right)
\Longleftrightarrow
\forall x\in E,\ \neg P(x).
\]
\end{theorembox}

\begin{preuvebox}
Dire que « tous les éléments de $E$ vérifient $P$ » est faux signifie précisément qu'il y a au moins un élément de $E$ qui ne vérifie pas $P$. Cela donne la première équivalence. Dire qu'« il existe un élément vérifiant $P$ » est faux signifie qu'aucun élément ne vérifie $P$, c'est-à-dire que tout élément vérifie $\neg P$. Cela donne la seconde équivalence.
\end{preuvebox}

\begin{exemplebox}{Négation complète}{neg-complete}
Nions la proposition :
\[
\forall \varepsilon>0,\exists \eta>0,\forall x\in\R,\ |x-a|<\eta\Rightarrow |f(x)-f(a)|<\varepsilon.
\]
On obtient :
\[
\exists \varepsilon>0,\forall \eta>0,\exists x\in\R,\ \neg\big(|x-a|<\eta\Rightarrow |f(x)-f(a)|<\varepsilon\big).
\]
Or la négation de $A\Rightarrow B$ est $A\land\neg B$. Donc :
\[
\exists \varepsilon>0,\forall \eta>0,\exists x\in\R,\ |x-a|<\eta\ \text{et}\ |f(x)-f(a)|\geq \varepsilon.
\]
\end{exemplebox}

\subsection{Exercices d'application et bilan}

\begin{exercicebox}{Lire les quantificateurs}{exo-lire-quant}
Dire si les propositions suivantes sont vraies ou fausses. Justifier.
\begin{enumerate}
  \item $\forall x\in\R,\exists y\in\R,\ x+y=0$.
  \item $\exists y\in\R,\forall x\in\R,\ x+y=0$.
  \item $\forall x\in\R,\exists y\in\R,\ y^2=x^2$.
  \item $\exists y\in\R,\forall x\in\R,\ y^2=x^2$.
\end{enumerate}
\end{exercicebox}

\begin{exercicebox}{Négations quantifiées}{exo-neg-quant}
Nier les propositions suivantes sans utiliser le symbole $\neg$ devant une proposition composée :
\begin{enumerate}
  \item $\forall n\in\N,\exists p\in\N,\ p>n$.
  \item $\exists M>0,\forall x\in\R,\ |f(x)|\leq M$.
  \item $\forall \varepsilon>0,\exists N\in\N,\forall n\geq N,\ |u_n-\ell|<\varepsilon$.
\end{enumerate}
\end{exercicebox}


%============================================================
\section{Rédiger en mathématiques : variables, quantificateurs et conventions}
%============================================================

\subsection{Pourquoi une partie sur la rédaction ?}

En mathématiques supérieures, il ne suffit pas d'avoir une idée correcte : il faut être capable de l'exprimer dans une langue précise. Une preuve mathématique est un texte argumenté dans lequel chaque objet introduit doit avoir un statut clair.

\begin{retenirbox}
Une rédaction mathématique correcte doit permettre de répondre à tout instant aux questions suivantes :
\begin{itemize}
    \item Quels sont les objets que l'on considère ?
    \item Sont-ils fixés une fois pour toutes, arbitraires, quelconques, ou existentiels ?
    \item De quelles variables dépendent-ils ?
    \item Quelle proposition cherche-t-on à démontrer ?
    \item Quels résultats déjà établis utilise-t-on ?
    \item Les quantificateurs sont-ils respectés dans le bon ordre ?
\end{itemize}
\end{retenirbox}

La majorité des erreurs de rédaction en début de prépa vient d'une mauvaise gestion des variables. Il est donc essentiel de comprendre le statut exact des lettres que l'on manipule.

\begin{remarquebox}
Faire des mathématiques, ce n'est pas seulement calculer. C'est produire un discours où chaque symbole a un rôle précis. Une lettre peut représenter un nombre fixé, une inconnue, une variable quelconque, un paramètre, une indéterminée, un indice ou encore une variable muette. Confondre ces rôles conduit très souvent à des raisonnements faux.
\end{remarquebox}

%------------------------------------------------------------
\subsection{Objets mathématiques et variables}
%------------------------------------------------------------

\begin{definitionbox}{Variable}{}
Une \emph{variable} est un symbole, souvent une lettre, utilisé pour désigner un objet mathématique susceptible de varier dans un ensemble donné.

Par exemple, dans l'expression
\[
x^2+1,
\]
la lettre $x$ est une variable réelle si l'on précise que $x\in\mathbb{R}$.
\end{definitionbox}

\begin{definitionbox}{Ensemble de variation}{}
L'\emph{ensemble de variation} d'une variable est l'ensemble dans lequel cette variable prend ses valeurs.

Par exemple, dans l'énoncé
\[
\forall n\in\mathbb{N}, \quad n^2+n \text{ est pair},
\]
la variable $n$ varie dans $\mathbb{N}$.
\end{definitionbox}

\begin{erreurbox}
Écrire une expression avec une lettre sans préciser où elle vit est souvent une faute de rédaction.

Par exemple, écrire seulement
\[
x^2\geq 0
\]
est incomplet si l'on ne sait pas si $x$ est réel, complexe, une matrice, une fonction, etc.

Une rédaction correcte serait :
\[
\forall x\in\mathbb{R},\quad x^2\geq 0.
\]
\end{erreurbox}

\begin{exemplebox}{Différents statuts possibles d'une même expression}{}
L'expression
\[
x^2-2x+1
\]
peut avoir plusieurs sens selon le statut de $x$ :
\begin{itemize}
    \item si $x\in\mathbb{R}$, c'est un nombre réel ;
    \item si $x\in\mathbb{C}$, c'est un nombre complexe ;
    \item si $x$ est une matrice carrée, c'est une matrice ;
    \item si $x$ est une fonction, cela peut désigner une fonction obtenue par opérations.
\end{itemize}

Ainsi, une même écriture formelle peut désigner des objets de nature différente.
\end{exemplebox}

%------------------------------------------------------------
\subsection{Variable libre et variable liée}
%------------------------------------------------------------

\begin{definitionbox}{Variable libre}{}
Une variable est dite \emph{libre} dans une expression ou une proposition si elle n'est pas placée sous la portée d'un quantificateur.

Par exemple, dans
\[
x^2\geq 0,
\]
la variable $x$ est libre.
\end{definitionbox}

\begin{definitionbox}{Variable liée}{}
Une variable est dite \emph{liée} lorsqu'elle est introduite par un quantificateur.

Par exemple, dans
\[
\forall x\in\mathbb{R},\quad x^2\geq 0,
\]
la variable $x$ est liée par le quantificateur $\forall x\in\mathbb{R}$.
\end{definitionbox}

\begin{exemplebox}{Variable libre et variable liée}{}
Dans la proposition
\[
\forall x\in\mathbb{R},\quad x^2+a x+1\geq 0,
\]
la variable $x$ est liée, mais la lettre $a$ est libre.

Cette proposition dépend donc du paramètre $a$.
\end{exemplebox}

\begin{retenirbox}
Une proposition avec une variable libre n'est pas encore une affirmation complètement déterminée. Elle devient une vraie proposition lorsque :
\begin{itemize}
    \item on fixe la variable ;
    \item ou on la quantifie ;
    \item ou on précise qu'elle joue le rôle d'un paramètre déjà fixé.
\end{itemize}
\end{retenirbox}

\begin{exemplebox}{Importance de la quantification}{}
L'énoncé
\[
x^2=4
\]
n'est pas une proposition universelle. Sa vérité dépend de $x$.

Si $x=2$, elle est vraie.  
Si $x=3$, elle est fausse.

En revanche,
\[
\exists x\in\mathbb{R},\quad x^2=4
\]
est une proposition vraie, tandis que
\[
\forall x\in\mathbb{R},\quad x^2=4
\]
est une proposition fausse.
\end{exemplebox}

%------------------------------------------------------------
\subsection{Portée d'un quantificateur}
%------------------------------------------------------------

\begin{definitionbox}{Portée d'un quantificateur}{}
La \emph{portée} d'un quantificateur est la partie de la proposition dans laquelle la variable qu'il introduit est liée.
\end{definitionbox}

\begin{exemplebox}{Portée d'un quantificateur}{}
Dans la proposition
\[
\forall x\in\mathbb{R},\quad x^2\geq 0 \quad \text{et} \quad y>0,
\]
le quantificateur $\forall x$ ne lie que la variable $x$. La variable $y$ reste libre.

Pour éviter toute ambiguïté, on peut écrire :
\[
\forall x\in\mathbb{R},\quad \left(x^2\geq 0 \text{ et } y>0\right).
\]
Mais même dans ce cas, $y$ reste libre si aucun quantificateur ne l'introduit.
\end{exemplebox}

\begin{methodebox}
Pour vérifier qu'une proposition est bien rédigée, on peut entourer mentalement la portée de chaque quantificateur.

Par exemple :
\[
\forall x\in E,\quad \exists y\in F,\quad P(x,y).
\]

Ici :
\begin{itemize}
    \item $x$ est fixé arbitrairement dans $E$ ;
    \item une fois $x$ fixé, il existe un $y\in F$ ;
    \item ce $y$ peut dépendre de $x$.
\end{itemize}
\end{methodebox}

%------------------------------------------------------------
\subsection{Variable muette}
%------------------------------------------------------------

\begin{definitionbox}{Variable muette}{}
Une variable liée par un quantificateur, une somme, un produit, une intégrale ou une définition locale est appelée \emph{variable muette} lorsqu'elle n'a pas d'importance en elle-même et peut être renommée sans changer le sens de l'expression.
\end{definitionbox}

\begin{exemplebox}{Renommer une variable muette}{}
Les deux propositions suivantes ont exactement le même sens :
\[
\forall x\in\mathbb{R},\quad x^2\geq 0
\]
et
\[
\forall t\in\mathbb{R},\quad t^2\geq 0.
\]

La lettre utilisée n'a aucune importance : $x$ et $t$ sont ici des variables muettes.
\end{exemplebox}

\begin{exemplebox}{Variable muette dans une somme}{}
Les deux expressions
\[
\sum_{k=1}^{n} k^2
\]
et
\[
\sum_{j=1}^{n} j^2
\]
désignent la même somme. L'indice $k$ ou $j$ est une variable muette.

En revanche, dans
\[
\sum_{k=1}^{n} k^2,
\]
la lettre $n$ n'est pas muette : elle est un paramètre de la somme.
\end{exemplebox}

\begin{erreurbox}
Il ne faut pas réutiliser une variable déjà utilisée pour un autre rôle.

Par exemple, l'écriture
\[
\sum_{n=1}^{n} n^2
\]
est incorrecte, car la lettre $n$ joue deux rôles incompatibles :
\begin{itemize}
    \item borne supérieure ;
    \item indice de sommation.
\end{itemize}

On doit écrire par exemple :
\[
\sum_{k=1}^{n} k^2.
\]
\end{erreurbox}

%------------------------------------------------------------
\subsection{Variable fixée, variable quelconque et variable arbitraire}
%------------------------------------------------------------

\begin{definitionbox}{Objet fixé}{}
Dire que l'on \emph{fixe} un objet signifie qu'on le choisit une fois pour toutes dans un ensemble donné, puis qu'on travaille avec cet objet dans la suite du raisonnement.
\end{definitionbox}

\begin{definitionbox}{Objet quelconque}{}
Un objet est dit \emph{quelconque} lorsqu'il est choisi arbitrairement dans un ensemble, sans propriété supplémentaire.

Démontrer une propriété pour un objet quelconque de $E$ revient à démontrer une propriété universelle sur $E$.
\end{definitionbox}

\begin{exemplebox}{Choisir un réel quelconque}{}
Pour démontrer :
\[
\forall x\in\mathbb{R},\quad x^2+1>0,
\]
on commence par écrire :

\[
\text{Soit } x\in\mathbb{R}.
\]

Cela signifie que $x$ est désormais fixé, mais sans aucune hypothèse particulière autre que $x\in\mathbb{R}$.
\end{exemplebox}

\begin{retenirbox}
Pour démontrer une proposition de la forme
\[
\forall x\in E,\quad P(x),
\]
on rédige généralement :
\[
\text{Soit } x\in E. \text{ Montrons que } P(x).
\]

Le mot \og soit \fg{} ne signifie pas que l'on choisit un exemple particulier. Il signifie que l'on considère un élément arbitraire de $E$.
\end{retenirbox}

\begin{exemplebox}{Preuve universelle}{}
Démontrons que :
\[
\forall n\in\mathbb{Z},\quad n^2+n \text{ est pair}.
\]

\textbf{Rédaction.}

Soit $n\in\mathbb{Z}$. On a
\[
n^2+n=n(n+1).
\]
Or deux entiers consécutifs sont de parités opposées. Donc l'un des deux entiers $n$ et $n+1$ est pair. Par conséquent, leur produit $n(n+1)$ est pair.

Ainsi, pour tout $n\in\mathbb{Z}$, l'entier $n^2+n$ est pair.
\end{exemplebox}

%------------------------------------------------------------
\subsection{Paramètres}
%------------------------------------------------------------

\begin{definitionbox}{Paramètre}{}
Un \emph{paramètre} est une variable que l'on considère comme fixée pendant une partie du raisonnement, mais dont la valeur peut influencer le résultat final.
\end{definitionbox}

\begin{exemplebox}{Paramètre dans une équation}{}
Dans l'équation
\[
x^2-2ax+1=0,
\]
on peut considérer $x$ comme l'inconnue et $a$ comme un paramètre réel.

On étudie alors les solutions en fonction de $a$.
\end{exemplebox}

\begin{methodebox}
Lorsque l'on travaille avec un paramètre, il faut toujours préciser son ensemble d'appartenance.

On écrit par exemple :
\[
\text{Soit } a\in\mathbb{R}. \text{ On cherche les solutions } x\in\mathbb{R} \text{ de l'équation } x^2-2ax+1=0.
\]

Ainsi :
\begin{itemize}
    \item $a$ est fixé ;
    \item $x$ est l'inconnue ;
    \item les solutions trouvées peuvent dépendre de $a$.
\end{itemize}
\end{methodebox}

\begin{exemplebox}{Solution dépendant d'un paramètre}{}
Étudions, selon le paramètre $a\in\mathbb{R}$, l'équation
\[
x+a=0.
\]

Pour tout $a\in\mathbb{R}$, l'unique solution est
\[
x=-a.
\]

On ne dit pas que $x$ est égal à un nombre fixe : on dit que la solution dépend du paramètre $a$.
\end{exemplebox}

%------------------------------------------------------------
\subsection{Dépendance entre variables}
%------------------------------------------------------------

La dépendance entre variables est un point fondamental.

\begin{exemplebox}{Ordre des quantificateurs et dépendance}{}
Les deux propositions suivantes sont très différentes :
\[
\forall x\in\mathbb{R},\quad \exists y\in\mathbb{R},\quad y>x
\]
et
\[
\exists y\in\mathbb{R},\quad \forall x\in\mathbb{R},\quad y>x.
\]

La première est vraie : pour chaque réel $x$, on peut prendre par exemple $y=x+1$.

La seconde est fausse : il n'existe pas de réel $y$ plus grand que tous les réels $x$.
\end{exemplebox}

Dans la proposition
\[
\forall x\in E,\quad \exists y\in F,\quad P(x,y),
\]
la variable $y$ peut dépendre de $x$.

Dans la proposition
\[
\exists y\in F,\quad \forall x\in E,\quad P(x,y),
\]
le même $y$ doit fonctionner pour tous les $x$.

\begin{erreurbox}
Confondre
\[
\forall x,\exists y
\]
et
\[
\exists y,\forall x
\]
est une erreur très grave.

L'ordre des quantificateurs change presque toujours le sens de l'énoncé.
\end{erreurbox}

\begin{exemplebox}{Majorant dépendant ou non de la variable}{}
Considérons la proposition :
\[
\forall x\in\mathbb{R},\quad \exists M\in\mathbb{R},\quad x\leq M.
\]
Elle est vraie : pour chaque $x$, on peut choisir $M=x$ ou $M=x+1$.

Mais la proposition :
\[
\exists M\in\mathbb{R},\quad \forall x\in\mathbb{R},\quad x\leq M
\]
est fausse : aucun réel $M$ ne majore tout $\mathbb{R}$.
\end{exemplebox}

%------------------------------------------------------------
\subsection{Constantes, inconnues, indéterminées}
%------------------------------------------------------------

\begin{definitionbox}{Constante}{}
Une \emph{constante} est un objet fixé dont la valeur ne varie pas dans le raisonnement.
\end{definitionbox}

\begin{definitionbox}{Inconnue}{}
Une \emph{inconnue} est une variable dont on cherche les valeurs satisfaisant une condition donnée.
\end{definitionbox}

\begin{definitionbox}{Indéterminée}{}
Une \emph{indéterminée} est une lettre utilisée formellement, par exemple dans un polynôme, sans nécessairement lui attribuer une valeur numérique.
\end{definitionbox}

\begin{exemplebox}{Inconnue et indéterminée}{}
Dans l'équation
\[
x^2-3x+2=0,
\]
$x$ est une inconnue.

Dans le polynôme
\[
P(X)=X^2-3X+2,
\]
$X$ est une indéterminée.

Dans l'expression
\[
P(a)=a^2-3a+2,
\]
$a$ désigne une valeur à laquelle on évalue le polynôme.
\end{exemplebox}

\begin{retenirbox}
En algèbre, on utilise souvent :
\begin{itemize}
    \item des minuscules $x,y,z$ pour des variables ou inconnues ;
    \item des majuscules $X,Y$ pour des indéterminées ;
    \item des lettres comme $a,b,\lambda,\mu$ pour des paramètres ou scalaires ;
    \item des lettres comme $n,p,k$ pour des entiers ou indices ;
    \item des lettres comme $E,F,V$ pour des ensembles ou espaces vectoriels.
\end{itemize}
Ces conventions ne sont pas absolues, mais elles aident à lire correctement un texte mathématique.
\end{retenirbox}

%------------------------------------------------------------
\subsection{Égalité, équation, identité}
%------------------------------------------------------------

\begin{definitionbox}{Égalité}{}
Une égalité est une affirmation de la forme
\[
A=B,
\]
où $A$ et $B$ désignent deux objets mathématiques de même nature.
\end{definitionbox}

\begin{definitionbox}{Équation}{}
Une équation est une égalité contenant une ou plusieurs inconnues, et que l'on cherche à résoudre.
\end{definitionbox}

\begin{definitionbox}{Identité}{}
Une identité est une égalité vraie pour toutes les valeurs admissibles des variables.
\end{definitionbox}

\begin{exemplebox}{Égalité, équation et identité}{}
L'égalité
\[
2+3=5
\]
est une égalité vraie.

L'écriture
\[
x^2-1=0
\]
est une équation si l'on cherche les valeurs de $x$ qui la vérifient.

L'égalité
\[
(x-1)(x+1)=x^2-1
\]
est une identité sur $\mathbb{R}$ :
\[
\forall x\in\mathbb{R},\quad (x-1)(x+1)=x^2-1.
\]
\end{exemplebox}

\begin{erreurbox}
Il ne faut pas confondre résoudre une équation et démontrer une identité.

Résoudre :
\[
x^2-1=0
\]
signifie chercher les $x$ tels que l'égalité soit vraie.

Démontrer :
\[
(x-1)(x+1)=x^2-1
\]
signifie prouver que l'égalité est vraie pour tout $x$ dans l'ensemble considéré.
\end{erreurbox}

%------------------------------------------------------------
\subsection{Bien introduire une variable dans une preuve}
%------------------------------------------------------------

Une variable ne doit jamais apparaître sans être introduite. Une rédaction correcte annonce toujours les objets utilisés.

\begin{methodebox}
Pour introduire correctement une variable, on utilise des phrases comme :
\begin{itemize}
    \item \og Soit $x\in E$ \fg{} ;
    \item \og Fixons $x\in E$ \fg{} ;
    \item \og Pour tout $x\in E$ \fg{} ;
    \item \og Il existe $x\in E$ tel que... \fg{} ;
    \item \og On cherche $x\in E$ vérifiant... \fg{} ;
    \item \og Soient $x,y\in E$ \fg{} ;
    \item \og Soit $\varepsilon>0$ \fg{} ;
    \item \og Soit $n\in\mathbb{N}$ \fg{}.
\end{itemize}
\end{methodebox}

\begin{exemplebox}{Rédaction incorrecte et rédaction correcte}{}
Rédaction incorrecte :
\[
x\in\mathbb{R},\quad x^2+1>0.
\]

Cette écriture n'est pas une phrase mathématique complète.

Rédaction correcte :
\[
\text{Soit } x\in\mathbb{R}. \text{ Alors } x^2\geq 0, \text{ donc } x^2+1\geq 1>0.
\]
\end{exemplebox}

%------------------------------------------------------------
\subsection{Les différents statuts possibles d'une variable}
%------------------------------------------------------------

Une même lettre peut avoir des statuts différents selon le contexte.

\begin{center}
\renewcommand{\arraystretch}{1.35}
\begin{tabularx}{\textwidth}{|>{\bfseries}l|X|X|}
\hline
Statut & Signification & Exemple \\
\hline
Variable libre & Lettre non quantifiée dont dépend l'expression & $x^2+1$ dépend de $x$ \\
\hline
Variable liée & Lettre introduite par un quantificateur & $\forall x\in\mathbb{R}, x^2\geq0$ \\
\hline
Variable muette & Variable liée que l'on peut renommer & $\sum_{k=1}^n k^2$ \\
\hline
Paramètre & Variable fixée dont dépend le résultat & $x^2-2ax+1=0$ \\
\hline
Inconnue & Variable que l'on cherche à déterminer & Résoudre $x^2=4$ \\
\hline
Constante & Objet fixé une fois pour toutes & $\pi$, $e$, ou un réel $a$ fixé \\
\hline
Indice & Variable entière servant à numéroter & $(u_n)_{n\in\mathbb{N}}$ \\
\hline
Indéterminée & Lettre formelle dans un polynôme & $P(X)=X^2+1$ \\
\hline
\end{tabularx}
\end{center}

%------------------------------------------------------------
\subsection{Rédiger une preuve universelle}
%------------------------------------------------------------

\begin{methodebox}
Pour démontrer :
\[
\forall x\in E,\quad P(x),
\]
on rédige :
\begin{enumerate}
    \item Soit $x\in E$.
    \item On démontre $P(x)$ sans supposer de propriété particulière sur $x$.
    \item On conclut : comme $x$ était quelconque dans $E$, la propriété est vraie pour tout $x\in E$.
\end{enumerate}
\end{methodebox}

\begin{exemplebox}{Preuve universelle rédigée}{}
Démontrons :
\[
\forall x\in\mathbb{R},\quad x^2-2x+2>0.
\]

\textbf{Rédaction.}

Soit $x\in\mathbb{R}$. On complète le carré :
\[
x^2-2x+2=(x-1)^2+1.
\]
Or, comme $x-1\in\mathbb{R}$, on a
\[
(x-1)^2\geq 0.
\]
Donc
\[
(x-1)^2+1\geq 1>0.
\]
Ainsi,
\[
x^2-2x+2>0.
\]
Comme $x$ était un réel quelconque, on a bien :
\[
\forall x\in\mathbb{R},\quad x^2-2x+2>0.
\]
\end{exemplebox}

%------------------------------------------------------------
\subsection{Rédiger une preuve existentielle}
%------------------------------------------------------------

\begin{methodebox}
Pour démontrer :
\[
\exists x\in E,\quad P(x),
\]
il suffit souvent de construire explicitement un exemple.

On rédige :
\begin{enumerate}
    \item On propose un candidat $x_0\in E$.
    \item On vérifie que $P(x_0)$ est vraie.
    \item On conclut qu'il existe bien un élément de $E$ vérifiant $P$.
\end{enumerate}
\end{methodebox}

\begin{exemplebox}{Preuve d'existence}{}
Démontrons :
\[
\exists x\in\mathbb{R},\quad x^2-3x+2=0.
\]

Prenons $x=1$. Alors
\[
1^2-3\cdot 1+2=1-3+2=0.
\]
Donc il existe bien un réel $x$ tel que
\[
x^2-3x+2=0.
\]
\end{exemplebox}

\begin{erreurbox}
Pour démontrer une existence, il ne suffit pas d'écrire :
\[
x^2-3x+2=0.
\]
Il faut soit donner un exemple, soit produire un argument prouvant qu'un tel objet existe.
\end{erreurbox}

%------------------------------------------------------------
\subsection{Rédiger avec plusieurs quantificateurs}
%------------------------------------------------------------

\begin{methodebox}
Pour rédiger une proposition avec plusieurs quantificateurs, il faut respecter leur ordre.

Pour démontrer :
\[
\forall x\in E,\quad \exists y\in F,\quad P(x,y),
\]
on écrit :
\begin{enumerate}
    \item Soit $x\in E$.
    \item On construit un $y\in F$, éventuellement en fonction de $x$.
    \item On vérifie $P(x,y)$.
\end{enumerate}
\end{methodebox}

\begin{exemplebox}{Rédaction avec deux quantificateurs}{}
Démontrons :
\[
\forall x\in\mathbb{R},\quad \exists y\in\mathbb{R},\quad y>x.
\]

Soit $x\in\mathbb{R}$. Posons
\[
y=x+1.
\]
Alors $y\in\mathbb{R}$ et
\[
y-x=(x+1)-x=1>0.
\]
Donc $y>x$.

Ainsi, pour tout réel $x$, il existe un réel $y$ tel que $y>x$.
\end{exemplebox}

\begin{retenirbox}
Dans l'exemple précédent, le réel $y$ dépend de $x$. On pourrait écrire :
\[
y=y(x)=x+1.
\]
Cette dépendance est autorisée parce que le quantificateur $\exists y$ vient après le quantificateur $\forall x$.
\end{retenirbox}

%------------------------------------------------------------
\subsection{Exemple complet : rédaction parfaite d'un exercice}
%------------------------------------------------------------

\begin{exercicebox}{Exemple de rédaction parfaite}{}
Soit $f:\mathbb{R}\to\mathbb{R}$ la fonction définie par
\[
f(x)=x^2-2x+3.
\]
\begin{enumerate}
    \item Montrer que, pour tout $x\in\mathbb{R}$, $f(x)\geq 2$.
    \item En déduire que $f$ admet un minimum sur $\mathbb{R}$.
    \item Déterminer l'ensemble des réels $x$ pour lesquels $f(x)=2$.
\end{enumerate}
\end{exercicebox}

\begin{preuvebox}[title={Correction détaillée et rédaction modèle}]

\textbf{Question 1.}

On veut démontrer une propriété universelle :
\[
\forall x\in\mathbb{R},\quad f(x)\geq 2.
\]

On commence donc par fixer un réel quelconque.

Soit $x\in\mathbb{R}$. Par définition de $f$, on a
\[
f(x)=x^2-2x+3.
\]
On complète le carré :
\[
x^2-2x+3=(x-1)^2+2.
\]
Or $x-1\in\mathbb{R}$, donc
\[
(x-1)^2\geq 0.
\]
Par conséquent,
\[
(x-1)^2+2\geq 2.
\]
Ainsi,
\[
f(x)\geq 2.
\]

Comme $x$ était un réel quelconque, on a bien démontré :
\[
\forall x\in\mathbb{R},\quad f(x)\geq 2.
\]

\medskip

\textbf{Question 2.}

La question précédente montre que $2$ est un minorant de l'ensemble des valeurs prises par $f$.

Pour montrer que $f$ admet un minimum égal à $2$, il faut maintenant vérifier que cette valeur est effectivement atteinte.

Calculons $f(1)$ :
\[
f(1)=1^2-2\cdot 1+3=1-2+3=2.
\]
Ainsi, la fonction $f$ prend la valeur $2$ en $x=1$.

D'après la question 1, pour tout $x\in\mathbb{R}$,
\[
f(x)\geq 2,
\]
et on vient de montrer que
\[
f(1)=2.
\]
Donc $f$ admet un minimum sur $\mathbb{R}$, et ce minimum vaut $2$.

\medskip

\textbf{Question 3.}

On cherche les réels $x$ tels que
\[
f(x)=2.
\]
Soit $x\in\mathbb{R}$. On a :
\[
f(x)=2
\iff x^2-2x+3=2
\]
\[
\iff x^2-2x+1=0
\]
\[
\iff (x-1)^2=0.
\]
Or, pour un réel $u$, on a
\[
u^2=0 \iff u=0.
\]
Donc
\[
(x-1)^2=0 \iff x-1=0 \iff x=1.
\]

Ainsi, l'ensemble des réels $x$ tels que $f(x)=2$ est
\[
\{1\}.
\]

\medskip

\textbf{Commentaires de rédaction.}

Dans cette correction :
\begin{itemize}
    \item la variable $x$ est introduite avant d'être utilisée ;
    \item on précise toujours dans quel ensemble elle vit ;
    \item la phrase \og Soit $x\in\mathbb{R}$ \fg{} signifie que $x$ est arbitraire ;
    \item la conclusion universelle est justifiée par le fait que $x$ était quelconque ;
    \item dans la question 3, $x$ n'est plus arbitraire : c'est une inconnue que l'on cherche à déterminer ;
    \item chaque équivalence est justifiée par une transformation algébrique correcte.
\end{itemize}

\end{preuvebox}

%------------------------------------------------------------
\subsection{Exercices d'application}
%------------------------------------------------------------

\begin{exercicebox}{Identifier les variables}{}
Dans chacune des expressions suivantes, identifier les variables libres, les variables liées et les paramètres éventuels.
\begin{enumerate}
    \item $\forall x\in\mathbb{R},\quad x^2+a\geq 0$.
    \item $\exists n\in\mathbb{N},\quad n>p$.
    \item $\sum_{k=0}^{n} k^2$.
    \item $\forall \varepsilon>0,\quad \exists \eta>0,\quad \forall x\in\mathbb{R},\quad |x-a|<\eta \Rightarrow |f(x)-f(a)|<\varepsilon$.
\end{enumerate}
\end{exercicebox}

\begin{exercicebox}{Réécriture correcte}{}
Réécrire correctement les phrases mathématiques suivantes.
\begin{enumerate}
    \item $x^2\geq 0$.
    \item $n^2+n$ est pair.
    \item $x^2-1=0$ admet deux solutions.
    \item Si $x>0$, alors il existe $y$ tel que $0<y<x$.
\end{enumerate}
\end{exercicebox}

\begin{exercicebox}{Comparer des quantifications}{}
Comparer les propositions suivantes et dire lesquelles sont vraies.
\begin{enumerate}
    \item $\forall x\in\mathbb{R},\quad \exists y\in\mathbb{R},\quad y>x$.
    \item $\exists y\in\mathbb{R},\quad \forall x\in\mathbb{R},\quad y>x$.
    \item $\forall x\in\mathbb{R},\quad \exists y\in\mathbb{R},\quad x+y=0$.
    \item $\exists y\in\mathbb{R},\quad \forall x\in\mathbb{R},\quad x+y=0$.
\end{enumerate}
\end{exercicebox}

\begin{exercicebox}{Équation avec paramètre}{}
Soit $a\in\mathbb{R}$. Résoudre, selon le paramètre $a$, l'équation d'inconnue $x\in\mathbb{R}$ :
\[
ax=1.
\]
\end{exercicebox}

\begin{exercicebox}{Rédaction complète}{}
Rédiger parfaitement la démonstration de la proposition suivante :
\[
\forall x\in\mathbb{R},\quad x^2+4x+5\geq 1.
\]
\end{exercicebox}

%------------------------------------------------------------
\subsection{Corrections des exercices}
%------------------------------------------------------------

\begin{preuvebox}[title={Correction de l'exercice 1}]

\begin{enumerate}
    \item Dans
    \[
    \forall x\in\mathbb{R},\quad x^2+a\geq 0,
    \]
    la variable $x$ est liée par le quantificateur $\forall x\in\mathbb{R}$. La lettre $a$ est libre : elle joue le rôle d'un paramètre.

    \item Dans
    \[
    \exists n\in\mathbb{N},\quad n>p,
    \]
    la variable $n$ est liée par le quantificateur existentiel. La lettre $p$ est libre : la vérité de la proposition dépend de $p$.

    \item Dans
    \[
    \sum_{k=0}^{n} k^2,
    \]
    la variable $k$ est une variable muette liée à la sommation. La lettre $n$ est un paramètre.

    \item Dans
    \[
    \forall \varepsilon>0,\quad \exists \eta>0,\quad \forall x\in\mathbb{R},\quad |x-a|<\eta \Rightarrow |f(x)-f(a)|<\varepsilon,
    \]
    les variables $\varepsilon$, $\eta$ et $x$ sont liées par des quantificateurs. Les objets $a$ et $f$ sont libres si le contexte ne les a pas déjà fixés.
\end{enumerate}

\end{preuvebox}

\begin{preuvebox}[title={Correction de l'exercice 2}]

\begin{enumerate}
    \item Une rédaction correcte est :
    \[
    \forall x\in\mathbb{R},\quad x^2\geq 0.
    \]

    \item Une rédaction correcte est :
    \[
    \forall n\in\mathbb{Z},\quad n^2+n \text{ est pair}.
    \]
    On peut aussi choisir $\mathbb{N}$ selon le contexte.

    \item Une rédaction correcte est :
    \[
    \text{L'équation } x^2-1=0 \text{ d'inconnue } x\in\mathbb{R} \text{ admet deux solutions.}
    \]
    On peut préciser :
    \[
    \{x\in\mathbb{R}\mid x^2-1=0\}=\{-1,1\}.
    \]

    \item Une rédaction correcte est :
    \[
    \forall x\in\mathbb{R}_+^*,\quad \exists y\in\mathbb{R},\quad 0<y<x.
    \]
    On peut même écrire :
    \[
    \forall x>0,\quad \exists y>0,\quad y<x.
    \]
\end{enumerate}

\end{preuvebox}

\begin{preuvebox}[title={Correction de l'exercice 3}]

\begin{enumerate}
    \item La proposition
    \[
    \forall x\in\mathbb{R},\quad \exists y\in\mathbb{R},\quad y>x
    \]
    est vraie. Pour un réel $x$ fixé, il suffit de choisir $y=x+1$.

    \item La proposition
    \[
    \exists y\in\mathbb{R},\quad \forall x\in\mathbb{R},\quad y>x
    \]
    est fausse. Elle signifierait qu'il existe un réel strictement plus grand que tous les réels, ce qui est impossible.

    \item La proposition
    \[
    \forall x\in\mathbb{R},\quad \exists y\in\mathbb{R},\quad x+y=0
    \]
    est vraie. Pour un réel $x$ fixé, on choisit $y=-x$.

    \item La proposition
    \[
    \exists y\in\mathbb{R},\quad \forall x\in\mathbb{R},\quad x+y=0
    \]
    est fausse. Elle demanderait un même réel $y$ tel que $x=-y$ pour tous les réels $x$, ce qui est impossible.
\end{enumerate}

\end{preuvebox}

\begin{preuvebox}[title={Correction de l'exercice 4}]

Soit $a\in\mathbb{R}$. On cherche les solutions $x\in\mathbb{R}$ de l'équation
\[
ax=1.
\]

On distingue deux cas.

\medskip

\textbf{Premier cas : $a=0$.}

L'équation devient
\[
0\cdot x=1,
\]
c'est-à-dire
\[
0=1,
\]
ce qui est impossible. Il n'y a donc aucune solution.

\medskip

\textbf{Deuxième cas : $a\neq 0$.}

On peut diviser par $a$. L'équation
\[
ax=1
\]
est équivalente à
\[
x=\frac{1}{a}.
\]
Il y a donc une unique solution :
\[
x=\frac{1}{a}.
\]

Ainsi :
\[
\boxed{
\begin{cases}
\text{si } a=0, \text{ l'équation n'a pas de solution ;}\\
\text{si } a\neq 0, \text{ l'équation admet l'unique solution } x=\dfrac{1}{a}.
\end{cases}
}
\]

\end{preuvebox}

\begin{preuvebox}[title={Correction de l'exercice 5}]

On veut démontrer :
\[
\forall x\in\mathbb{R},\quad x^2+4x+5\geq 1.
\]

Soit $x\in\mathbb{R}$. On complète le carré :
\[
x^2+4x+5=(x+2)^2+1.
\]
Or $x+2\in\mathbb{R}$, donc
\[
(x+2)^2\geq 0.
\]
Par conséquent,
\[
(x+2)^2+1\geq 1.
\]
Ainsi,
\[
x^2+4x+5\geq 1.
\]

Comme $x$ était un réel quelconque, on a bien démontré :
\[
\forall x\in\mathbb{R},\quad x^2+4x+5\geq 1.
\]

\end{preuvebox}

%------------------------------------------------------------
\subsection{Bilan : règles pratiques de rédaction}
%------------------------------------------------------------

\begin{retenirbox}
\begin{enumerate}
    \item Ne jamais utiliser une lettre sans l'avoir introduite.
    \item Toujours préciser l'ensemble d'appartenance d'une variable.
    \item Pour démontrer un \og pour tout \fg{}, commencer par \og Soit $x\in E$ \fg{}.
    \item Pour démontrer une existence, construire un exemple ou justifier l'existence.
    \item Respecter l'ordre des quantificateurs.
    \item Ne pas confondre paramètre et inconnue.
    \item Ne pas réutiliser une variable muette comme variable libre.
    \item Vérifier les dépendances : un objet introduit après un quantificateur peut dépendre des objets introduits avant.
    \item Écrire des phrases, pas seulement des calculs.
    \item Terminer par une conclusion correspondant exactement à l'énoncé.
\end{enumerate}
\end{retenirbox}
%===========================================================
\section{Les grands types de raisonnement}

\subsection{Raisonnement direct}

\begin{methodebox}
Pour démontrer directement $P\Rightarrow Q$ :
\begin{enumerate}
  \item On suppose $P$ vraie.
  \item On utilise les définitions, les hypothèses et les théorèmes connus.
  \item On aboutit à $Q$.
\end{enumerate}
\end{methodebox}

\begin{exemplebox}{Raisonnement direct}{direct}
Montrons que, pour tout entier $n$, si $n$ est pair, alors $n^2$ est pair.

Soit $n\in\Z$. Supposons $n$ pair. Alors il existe $k\in\Z$ tel que $n=2k$. Donc
\[
n^2=(2k)^2=4k^2=2(2k^2).
\]
Comme $2k^2\in\Z$, on en déduit que $n^2$ est pair.
\end{exemplebox}

\subsection{Raisonnement par contraposée}

\begin{methodebox}
Pour démontrer $P\Rightarrow Q$ par contraposée, on démontre $\neg Q\Rightarrow \neg P$. Cette méthode est particulièrement efficace lorsque la négation de la conclusion est plus exploitable que l'hypothèse de départ.
\end{methodebox}

\begin{exemplebox}{Contraposée}{contraposee-ex}
Montrons que, pour tout $n\in\Z$, si $n^2$ est pair, alors $n$ est pair.

La contraposée est : si $n$ n'est pas pair, alors $n^2$ n'est pas pair. Supposons donc $n$ impair. Il existe $k\in\Z$ tel que $n=2k+1$. Alors
\[
n^2=(2k+1)^2=4k^2+4k+1=2(2k^2+2k)+1.
\]
Donc $n^2$ est impair. La contraposée est démontrée ; l'implication initiale est donc vraie.
\end{exemplebox}

\subsection{Raisonnement par l'absurde}

\begin{methodebox}
Pour démontrer une proposition $P$ par l'absurde :
\begin{enumerate}
  \item On suppose $\neg P$.
  \item On déduit une contradiction, c'est-à-dire une proposition manifestement fausse ou incompatible avec une hypothèse.
  \item On conclut que $\neg P$ est fausse, donc que $P$ est vraie.
\end{enumerate}
\end{methodebox}

\begin{exemplebox}{Irrationalité de $\sqrt2$}{sqrt2}
Montrons que $\sqrt2$ est irrationnel. Supposons par l'absurde que $\sqrt2\in\Q$. Alors il existe deux entiers $p,q\in\Z$, avec $q\neq0$, premiers entre eux, tels que
\[
\sqrt2=\frac pq.
\]
Donc $p^2=2q^2$. Ainsi $p^2$ est pair, donc $p$ est pair. Il existe donc $r\in\Z$ tel que $p=2r$. Alors
\[
4r^2=2q^2,
\]
donc $q^2=2r^2$. Ainsi $q^2$ est pair, donc $q$ est pair. Les entiers $p$ et $q$ sont donc tous deux pairs, contradiction avec le fait qu'ils sont premiers entre eux. Donc $\sqrt2$ est irrationnel.
\end{exemplebox}

\subsection{Raisonnement par disjonction de cas}

\begin{methodebox}
Pour prouver $P$ par disjonction de cas, on découpe la situation en cas exhaustifs et incompatibles ou non. Il faut vérifier deux choses :
\begin{itemize}
  \item les cas couvrent toutes les possibilités ;
  \item dans chacun des cas, la conclusion est démontrée.
\end{itemize}
\end{methodebox}

\begin{exemplebox}{Cas selon le signe}{cas-signe}
Montrons que pour tout $x\in\R$, $|x|\geq0$. Si $x\geq0$, alors $|x|=x\geq0$. Si $x<0$, alors $|x|=-x>0$. Dans tous les cas, $|x|\geq0$.
\end{exemplebox}

\subsection{Raisonnement par équivalences}

\begin{methodebox}
Une chaîne d'équivalences doit être utilisée avec prudence. Chaque passage doit être réversible. Si une étape n'est valable que dans un sens, il faut utiliser une implication et non le symbole $\Longleftrightarrow$.
\end{methodebox}

\begin{exemplebox}{Résolution par équivalences}{equiv-ex}
Résolvons dans $\R$ l'équation $|x-1|<2$.
\[
|x-1|<2 \Longleftrightarrow -2<x-1<2 \Longleftrightarrow -1<x<3.
\]
L'ensemble solution est donc $]-1,3[$.
\end{exemplebox}

\subsection{Raisonnement par analyse-synthèse}

\begin{methodebox}
Le raisonnement par analyse-synthèse est utile pour trouver un objet.
\begin{enumerate}
  \item \textbf{Analyse :} on suppose que l'objet existe et on déduit les conditions qu'il doit satisfaire. Cette étape guide la recherche.
  \item \textbf{Synthèse :} on vérifie que l'objet obtenu convient effectivement.
\end{enumerate}
\end{methodebox}

\begin{exemplebox}{Décomposition paire-impaire}{analyse-synthese}
Soit $f:\R\to\R$. On cherche à écrire $f$ sous la forme $f=p+i$, où $p$ est paire et $i$ impaire.

\textbf{Analyse.} Si $f=p+i$, avec $p$ paire et $i$ impaire, alors
\[
f(-x)=p(-x)+i(-x)=p(x)-i(x).
\]
En additionnant et en soustrayant :
\[
p(x)=\frac{f(x)+f(-x)}2,
\qquad
 i(x)=\frac{f(x)-f(-x)}2.
\]

\textbf{Synthèse.} Définissons
\[
p(x)=\frac{f(x)+f(-x)}2,
\qquad
 i(x)=\frac{f(x)-f(-x)}2.
\]
Alors $p$ est paire, $i$ est impaire, et $p+i=f$. La décomposition existe et elle est unique par les formules obtenues dans l'analyse.
\end{exemplebox}

\subsection{Raisonnement par récurrence}

\begin{theorembox}{Principe de récurrence simple}{rec-simple}
Soit $P(n)$ une propriété dépendant de $n\in\N$. Si :
\begin{enumerate}
  \item $P(0)$ est vraie ;
  \item pour tout $n\in\N$, $P(n)\Rightarrow P(n+1)$ ;
\end{enumerate}
alors $P(n)$ est vraie pour tout $n\in\N$.
\end{theorembox}

\begin{methodebox}
Une preuve par récurrence doit comporter explicitement :
\begin{enumerate}
  \item l'énoncé clair de la propriété $P(n)$ ;
  \item l'initialisation ;
  \item l'hérédité : on fixe $n$, on suppose $P(n)$ vraie, puis on démontre $P(n+1)$ ;
  \item la conclusion par le principe de récurrence.
\end{enumerate}
\end{methodebox}

\begin{exemplebox}{Somme des premiers entiers}{somme-entiers}
Montrons que pour tout $n\in\N$,
\[
\sum_{k=0}^n k=\frac{n(n+1)}2.
\]
Soit $P(n)$ cette propriété.

\textbf{Initialisation.} Pour $n=0$, on a $\sum_{k=0}^0 k=0$ et $0(0+1)/2=0$. Donc $P(0)$ est vraie.

\textbf{Hérédité.} Soit $n\in\N$. Supposons $P(n)$ vraie. Alors
\[
\sum_{k=0}^{n+1}k=\sum_{k=0}^{n}k+(n+1)
=\frac{n(n+1)}2+(n+1)
=\frac{(n+1)(n+2)}2.
\]
Donc $P(n+1)$ est vraie.

Par récurrence, la formule est vraie pour tout $n\in\N$.
\end{exemplebox}

\begin{theorembox}{Récurrence forte}{rec-forte}
Soit $P(n)$ une propriété dépendant de $n\in\N$. Si :
\begin{enumerate}
  \item $P(0)$ est vraie ;
  \item pour tout $n\in\N$, si $P(0),P(1),\dots,P(n)$ sont vraies, alors $P(n+1)$ est vraie ;
\end{enumerate}
alors $P(n)$ est vraie pour tout $n\in\N$.
\end{theorembox}

\begin{exemplebox}{Décomposition en facteurs premiers, version existence}{facteurs-premiers}
Montrons que tout entier $n\geq2$ admet une décomposition comme produit de nombres premiers.

On raisonne par récurrence forte sur $n\geq2$. Pour $n=2$, c'est vrai car $2$ est premier. Supposons le résultat vrai pour tous les entiers $k$ tels que $2\leq k\leq n$. Considérons $n+1$. Si $n+1$ est premier, le résultat est vrai. Sinon, il existe $a,b\in\N$ tels que $n+1=ab$, avec $2\leq a\leq n$ et $2\leq b\leq n$. Par hypothèse de récurrence forte, $a$ et $b$ sont produits de nombres premiers. Donc $n+1=ab$ l'est aussi. La propriété est démontrée.
\end{exemplebox}

\subsection{Contre-exemple}

\begin{methodebox}
Pour réfuter une proposition universelle $\forall x\in E, P(x)$, il suffit de trouver un élément $a\in E$ tel que $P(a)$ soit fausse. Cet élément s'appelle un contre-exemple.
\end{methodebox}

\begin{exemplebox}{Contre-exemple}{contre-ex}
La proposition « pour tout réel $x$, $x^2\geq x$ » est fausse. En effet, pour $x=\frac12$, on a $x^2=\frac14<\frac12=x$.
\end{exemplebox}

\subsection{Exercices d'application et bilan}

\begin{exercicebox}{Choisir la méthode}{exo-choix-methode}
Pour chacune des propositions suivantes, indiquer une méthode adaptée, puis rédiger une preuve complète.
\begin{enumerate}
  \item Pour tout $n\in\Z$, si $n^2$ est impair, alors $n$ est impair.
  \item Pour tout $n\in\N$, $1+2+\cdots+n=\dfrac{n(n+1)}2$.
  \item Il n'existe pas de plus grand entier naturel.
  \item Pour tout réel $x$, $x^2+2x+2>0$.
\end{enumerate}
\end{exercicebox}

\begin{exercicebox}{Récurrence plus difficile}{exo-rec-diff}
Montrer que pour tout $n\geq1$,
\[
\sum_{k=1}^n k^2=\frac{n(n+1)(2n+1)}6.
\]
\end{exercicebox}

%===========================================================
\section{Ensembles, appartenance, inclusion et opérations}

\subsection{Ensembles et éléments}

\begin{definitionbox}{Ensemble et appartenance}{ensemble}
Un ensemble est une collection d'objets appelés éléments. Si $x$ est un élément de $E$, on écrit $x\in E$. Sinon, on écrit $x\notin E$.
\end{definitionbox}

\begin{definitionbox}{Inclusion}{inclusion}
Soient $A$ et $B$ deux ensembles. On dit que $A$ est inclus dans $B$, et on écrit $A\subset B$ ou $A\subseteq B$, lorsque
\[
\forall x,\quad x\in A\Rightarrow x\in B.
\]
\end{definitionbox}

\begin{methodebox}
Pour montrer $A\subset B$, on prend un élément arbitraire $x\in A$, puis on démontre que $x\in B$.

Pour montrer $A=B$, on démontre souvent deux inclusions : $A\subset B$ et $B\subset A$.
\end{methodebox}

\subsection{Opérations ensemblistes}

\begin{definitionbox}{Réunion, intersection, complémentaire}{operations-ens}
Soient $A$ et $B$ deux parties d'un ensemble $E$.
\begin{align*}
A\cap B&=\{x\in E\mid x\in A\text{ et }x\in B\},\\
A\cup B&=\{x\in E\mid x\in A\text{ ou }x\in B\},\\
E\setminus A&=\{x\in E\mid x\notin A\}.
\end{align*}
\end{definitionbox}

\begin{propositionbox}{Lien logique-ensembles}{logique-ensembles}
Les opérations ensemblistes traduisent les connecteurs logiques :
\[
x\in A\cap B \Longleftrightarrow (x\in A)\land(x\in B),
\]
\[
x\in A\cup B \Longleftrightarrow (x\in A)\lor(x\in B),
\]
\[
x\in E\setminus A \Longleftrightarrow \neg(x\in A).
\]
\end{propositionbox}

\begin{theorembox}{Lois de De Morgan ensemblistes}{de-morgan-ens}
Soient $A$ et $B$ deux parties d'un ensemble $E$. Alors :
\[
E\setminus(A\cap B)=(E\setminus A)\cup(E\setminus B),
\]
\[
E\setminus(A\cup B)=(E\setminus A)\cap(E\setminus B).
\]
\end{theorembox}

\begin{preuvebox}
Montrons la première égalité par double inclusion, ou plus directement par équivalences. Pour tout $x\in E$,
\begin{align*}
x\in E\setminus(A\cap B)
&\Longleftrightarrow x\notin A\cap B\\
&\Longleftrightarrow \neg(x\in A\text{ et }x\in B)\\
&\Longleftrightarrow x\notin A\text{ ou }x\notin B\\
&\Longleftrightarrow x\in(E\setminus A)\cup(E\setminus B).
\end{align*}
Donc les deux ensembles ont les mêmes éléments, ils sont égaux.
\end{preuvebox}

\subsection{Produit cartésien et parties}

\begin{definitionbox}{Produit cartésien}{produit-cartesien}
Le produit cartésien $E\times F$ est l'ensemble des couples $(x,y)$ tels que $x\in E$ et $y\in F$ :
\[
E\times F=\{(x,y)\mid x\in E,\ y\in F\}.
\]
\end{definitionbox}

\begin{definitionbox}{Ensemble des parties}{parties}
L'ensemble des parties de $E$, noté $\calP(E)$, est l'ensemble dont les éléments sont les sous-ensembles de $E$.
\end{definitionbox}

\begin{exemplebox}{Parties d'un ensemble fini}{parties-ex}
Si $E=\{a,b\}$, alors
\[
\calP(E)=\{\varnothing,\{a\},\{b\},\{a,b\}\}.
\]
Il faut bien distinguer $a$, qui est un élément de $E$, et $\{a\}$, qui est une partie de $E$.
\end{exemplebox}

\begin{erreurbox}
On écrit $a\in E$, mais $\{a\}\subset E$. Écrire $\{a\}\in E$ est généralement faux, sauf si $E$ contient lui-même un ensemble comme élément.
\end{erreurbox}

\subsection{Exercices d'application et bilan}

\begin{exercicebox}{Démonstration ensembliste}{exo-ens1}
Soient $A,B,C$ trois parties d'un ensemble $E$. Montrer que
\[
A\cap(B\cup C)=(A\cap B)\cup(A\cap C).
\]
\end{exercicebox}

\begin{exercicebox}{Double inclusion}{exo-double-inclusion}
Soient $A$ et $B$ deux parties de $E$. Montrer que
\[
A\subset B \Longleftrightarrow A\cap B=A.
\]
\end{exercicebox}

%===========================================================
\section{Applications, images, antécédents et raisonnements associés}

\subsection{Définition d'une application}

\begin{definitionbox}{Application}{application}
Une application $f:E\to F$ associe à chaque élément $x\in E$ un unique élément de $F$, noté $f(x)$. L'ensemble $E$ est l'ensemble de départ, $F$ l'ensemble d'arrivée.
\end{definitionbox}

\begin{definitionbox}{Image directe et image réciproque}{image-reciproque}
Soit $f:E\to F$.
\begin{itemize}
  \item Si $A\subset E$, l'image directe de $A$ est
  \[
  f(A)=\{f(x)\mid x\in A\}.
  \]
  \item Si $B\subset F$, l'image réciproque de $B$ est
  \[
  f^{-1}(B)=\{x\in E\mid f(x)\in B\}.
  \]
\end{itemize}
Attention : $f^{-1}(B)$ existe toujours comme ensemble, même si $f$ n'est pas bijective. Il ne faut pas le confondre avec une fonction réciproque.
\end{definitionbox}

\subsection{Injection, surjection, bijection}

\begin{definitionbox}{Injection, surjection, bijection}{inj-surj-bij}
Soit $f:E\to F$.
\begin{itemize}
  \item $f$ est \emph{injective} si
  \[
  \forall x,x'\in E,\quad f(x)=f(x')\Rightarrow x=x'.
  \]
  \item $f$ est \emph{surjective} si
  \[
  \forall y\in F,\exists x\in E,\quad f(x)=y.
  \]
  \item $f$ est \emph{bijective} si elle est injective et surjective.
\end{itemize}
\end{definitionbox}

\begin{methodebox}
Pour prouver qu'une fonction est injective, on part de $f(x)=f(x')$ et on cherche à démontrer $x=x'$.

Pour prouver qu'elle est surjective, on fixe un $y$ quelconque dans l'ensemble d'arrivée et on cherche un antécédent $x$ tel que $f(x)=y$.
\end{methodebox}

\begin{exemplebox}{Injection et surjection}{inj-surj-ex}
Considérons $f:\R\to\R$, $f(x)=2x+1$.

\textbf{Injection.} Si $f(x)=f(x')$, alors $2x+1=2x'+1$, donc $x=x'$. Ainsi $f$ est injective.

\textbf{Surjection.} Soit $y\in\R$. Cherchons $x\in\R$ tel que $2x+1=y$. On obtient $x=(y-1)/2$, qui est réel. Donc tout $y\in\R$ a un antécédent. Ainsi $f$ est surjective.

La fonction $f$ est donc bijective.
\end{exemplebox}

\begin{propositionbox}{Image réciproque et opérations}{image-reciproque-operations}
Soit $f:E\to F$ et soient $B,C\subset F$. Alors :
\[
f^{-1}(B\cap C)=f^{-1}(B)\cap f^{-1}(C),
\]
\[
f^{-1}(B\cup C)=f^{-1}(B)\cup f^{-1}(C),
\]
\[
f^{-1}(F\setminus B)=E\setminus f^{-1}(B).
\]
\end{propositionbox}

\begin{preuvebox}
Montrons la première égalité. Pour tout $x\in E$,
\begin{align*}
x\in f^{-1}(B\cap C)
&\Longleftrightarrow f(x)\in B\cap C\\
&\Longleftrightarrow f(x)\in B\text{ et }f(x)\in C\\
&\Longleftrightarrow x\in f^{-1}(B)\text{ et }x\in f^{-1}(C)\\
&\Longleftrightarrow x\in f^{-1}(B)\cap f^{-1}(C).
\end{align*}
Donc les deux ensembles sont égaux.
\end{preuvebox}

\subsection{Exercices d'application et bilan}

\begin{exercicebox}{Injectivité et surjectivité}{exo-inj-surj}
Étudier l'injectivité, la surjectivité et la bijectivité des applications suivantes :
\begin{enumerate}
  \item $f:\R\to\R$, $f(x)=x^2$ ;
  \item $g:\R_+\to\R_+$, $g(x)=x^2$ ;
  \item $h:\R\to\R$, $h(x)=x^3+x$.
\end{enumerate}
\end{exercicebox}

\begin{exercicebox}{Images et images réciproques}{exo-images}
Soit $f:E\to F$ et soient $A,B\subset E$. Montrer que
\[
f(A\cup B)=f(A)\cup f(B).
\]
Montrer ensuite que l'inclusion
\[
f(A\cap B)\subset f(A)\cap f(B)
\]
est toujours vraie, puis donner un exemple où l'égalité est fausse.
\end{exercicebox}

%===========================================================
\section{Rédaction mathématique : de l'idée à la preuve}

\subsection{Le statut des objets}

Une preuve rigoureuse doit toujours préciser le statut des objets manipulés : sont-ils fixés ? arbitraires ? construits ? supposés exister ? dépendants d'un autre objet ?

\begin{methodebox}
Quelques phrases standard utiles :
\begin{itemize}
  \item « Soit $x\in E$ » : on fixe un élément arbitraire de $E$.
  \item « Supposons que... » : on introduit une hypothèse locale.
  \item « Il existe donc... » : on utilise une existence démontrée ou une définition.
  \item « Posons... » : on construit explicitement un objet.
  \item « Réciproquement... » : on commence l'autre implication.
  \item « D'où... » : on annonce une conséquence logique.
\end{itemize}
\end{methodebox}

\subsection{Passer d'une figure à une preuve}

Une figure peut suggérer une propriété, mais elle ne la démontre pas. En géométrie comme en analyse, il faut traduire l'information visuelle en énoncés mathématiques.

\begin{methodebox}
Pour passer d'une figure à une preuve :
\begin{enumerate}
  \item Identifier ce qui est donné par l'énoncé et ce qui est seulement observé sur la figure.
  \item Traduire les données en relations mathématiques : égalités de longueurs, angles, parallélisme, coordonnées, appartenance.
  \item Choisir un théorème utilisable avec ces hypothèses.
  \item Vérifier explicitement les conditions d'application du théorème.
  \item Conclure sans invoquer l'apparence de la figure.
\end{enumerate}
\end{methodebox}

\begin{exemplebox}{Figure et preuve}{figure-preuve}
Si une figure semble montrer que deux droites sont parallèles, cela ne suffit pas. Pour le prouver, on peut par exemple montrer que deux angles alternes-internes sont égaux, ou que les droites ont le même vecteur directeur, ou utiliser la réciproque du théorème de Thalès selon le contexte.
\end{exemplebox}

\subsection{Les symboles ne remplacent pas la pensée}

\begin{erreurbox}
Une suite de symboles sans phrases n'est pas toujours une démonstration. Les phrases servent à indiquer le statut des variables, le sens des implications, les hypothèses utilisées et la conclusion visée.
\end{erreurbox}

\begin{retenirbox}
Une bonne preuve est lisible par quelqu'un qui ne connaît pas votre brouillon. Elle dit clairement ce que l'on suppose, ce que l'on veut montrer, et pourquoi chaque étape est vraie.
\end{retenirbox}

\subsection{Exercices de rédaction}

\begin{exercicebox}{Corriger une preuve}{exo-corriger-preuve}
Un élève écrit :
\[
x^2=1 \Longleftrightarrow x=1.
\]
Expliquer l'erreur, puis rédiger correctement la résolution de $x^2=1$ dans $\R$.
\end{exercicebox}

\begin{exercicebox}{Rédiger une double implication}{exo-redaction-equiv}
Démontrer soigneusement que, pour tout entier $n\in\Z$,
\[
n\text{ est pair }\Longleftrightarrow n^2\text{ est pair}.
\]
\end{exercicebox}

%===========================================================
\section{Corrections détaillées des exercices}

\subsection*{Correction des exercices de la section 1}
\addcontentsline{toc}{subsection}{Correction des exercices de la section 1}

\paragraph{Exercice \ref{exo:exo-prop}.}
\begin{enumerate}
  \item $3^2=9$ est une proposition vraie.
  \item $x+1=0$ est un prédicat, avec variable libre $x$.
  \item La phrase est une proposition vraie : pour tout $n\in\N$, $n^2+n=n(n+1)$ est produit de deux entiers consécutifs, donc pair.
  \item C'est une proposition fausse dans $\R$, car $x^2\geq0$ donc $x^2+1\geq1$.
  \item $x<y$ est un prédicat à deux variables libres $x$ et $y$.
\end{enumerate}

\paragraph{Exercice \ref{exo:exo-transformer}.}
Pour $x^2\geq0$, une proposition vraie est $\forall x\in\R, x^2\geq0$ ; une proposition fausse est $\forall x\in\C, x^2\geq0$, qui n'a même pas de sens avec l'ordre usuel sur $\C$, ou plus simplement $\forall x\in\R, x^2<0$. Pour $x^2=4$, une proposition vraie est $\exists x\in\R,x^2=4$ ; une proposition fausse est $\forall x\in\R,x^2=4$. Pour « $n$ est divisible par $2$ », une proposition vraie est $\exists n\in\N, n$ est divisible par $2$ ; une proposition fausse est $\forall n\in\N, n$ est divisible par $2$.

\subsection*{Correction des exercices de la section 2}
\addcontentsline{toc}{subsection}{Correction des exercices de la section 2}

\paragraph{Exercice \ref{exo:exo-table-verite}.}
Les propositions $P\Rightarrow(P\lor Q)$ et $(P\land Q)\Rightarrow P$ sont des tautologies, car dès que l'hypothèse est vraie, la conclusion l'est. La proposition $(P\Rightarrow Q)\Longleftrightarrow(\neg P\lor Q)$ est une tautologie : elle donne une forme équivalente de l'implication. En revanche, $(P\Rightarrow Q)\Longleftrightarrow(Q\Rightarrow P)$ n'est pas une tautologie : prendre $P$ vraie et $Q$ fausse donne $P\Rightarrow Q$ fausse, tandis que $Q\Rightarrow P$ est vraie.

\paragraph{Exercice \ref{exo:exo-negations}.}
\begin{enumerate}
  \item La négation de « $x>0$ et $x<1$ » est « $x\leq0$ ou $x\geq1$ ».
  \item La négation est : « $n$ n'est pas pair et $n$ n'est pas divisible par $3$ ».
  \item La négation de $x>2\Rightarrow x^2>4$ est $x>2$ et $x^2\leq4$.
  \item $P\Longleftrightarrow Q$ signifie même valeur de vérité. Sa négation est : $P$ et $Q$ ont des valeurs de vérité différentes, c'est-à-dire $(P\land\neg Q)\lor(Q\land\neg P)$.
\end{enumerate}

\subsection*{Correction des exercices de la section 3}
\addcontentsline{toc}{subsection}{Correction des exercices de la section 3}

\paragraph{Exercice \ref{exo:exo-lire-quant}.}
\begin{enumerate}
  \item Vraie : pour tout $x$, on prend $y=-x$.
  \item Fausse : il n'existe pas un seul réel $y$ tel que $x+y=0$ pour tous les réels $x$.
  \item Vraie : pour tout $x$, on peut prendre $y=x$.
  \item Fausse : si $y^2=x^2$ pour tout $x$, en particulier pour $x=0$ on a $y=0$, puis pour $x=1$ il faudrait $0=1$, contradiction.
\end{enumerate}

\paragraph{Exercice \ref{exo:exo-neg-quant}.}
\begin{enumerate}
  \item $\exists n\in\N,\forall p\in\N,\ p\leq n$.
  \item $\forall M>0,\exists x\in\R,\ |f(x)|>M$.
  \item $\exists \varepsilon>0,\forall N\in\N,\exists n\geq N,\ |u_n-\ell|\geq\varepsilon$.
\end{enumerate}

\subsection*{Correction des exercices de la section 4}
\addcontentsline{toc}{subsection}{Correction des exercices de la section 4}

\paragraph{Exercice \ref{exo:exo-choix-methode}.}
\begin{enumerate}
  \item On peut raisonner par contraposée. Si $n$ est pair, alors $n=2k$, donc $n^2=4k^2=2(2k^2)$ est pair. La contraposée de « $n^2$ impair implique $n$ impair » est donc démontrée.
  \item Récurrence simple, comme dans l'exemple du cours.
  \item Par l'absurde : supposons qu'il existe un plus grand entier naturel $N$. Alors $N+1\in\N$ et $N+1>N$, contradiction.
  \item Raisonnement direct par forme canonique : $x^2+2x+2=(x+1)^2+1>0$ pour tout réel $x$.
\end{enumerate}

\paragraph{Exercice \ref{exo:exo-rec-diff}.}
Soit $P(n)$ la propriété
\[
\sum_{k=1}^n k^2=\frac{n(n+1)(2n+1)}6.
\]
Initialisation : pour $n=1$, les deux membres valent $1$.
Supposons $P(n)$ vraie. Alors
\begin{align*}
\sum_{k=1}^{n+1}k^2
&=\sum_{k=1}^n k^2+(n+1)^2\\
&=\frac{n(n+1)(2n+1)}6+(n+1)^2\\
&=(n+1)\left(\frac{n(2n+1)}6+n+1\right)\\
&=(n+1)\frac{2n^2+n+6n+6}{6}\\
&=(n+1)\frac{2n^2+7n+6}{6}\\
&=(n+1)\frac{(n+2)(2n+3)}6\\
&=\frac{(n+1)(n+2)(2(n+1)+1)}6.
\end{align*}
Donc $P(n+1)$ est vraie. Par récurrence, la formule est vraie pour tout $n\geq1$.

\subsection*{Correction des exercices de la section 5}
\addcontentsline{toc}{subsection}{Correction des exercices de la section 5}

\paragraph{Exercice \ref{exo:exo-ens1}.}
Pour tout $x\in E$,
\begin{align*}
x\in A\cap(B\cup C)
&\Longleftrightarrow x\in A\text{ et }x\in B\cup C\\
&\Longleftrightarrow x\in A\text{ et }(x\in B\text{ ou }x\in C)\\
&\Longleftrightarrow (x\in A\text{ et }x\in B)\text{ ou }(x\in A\text{ et }x\in C)\\
&\Longleftrightarrow x\in(A\cap B)\cup(A\cap C).
\end{align*}
Donc les deux ensembles sont égaux.

\paragraph{Exercice \ref{exo:exo-double-inclusion}.}
Supposons $A\subset B$. Alors si $x\in A\cap B$, on a $x\in A$, donc $A\cap B\subset A$. Réciproquement, si $x\in A$, alors comme $A\subset B$, $x\in B$, donc $x\in A\cap B$. Ainsi $A\subset A\cap B$. Donc $A\cap B=A$.

Réciproquement, supposons $A\cap B=A$. Soit $x\in A$. Comme $A=A\cap B$, on a $x\in A\cap B$, donc $x\in B$. Ainsi $A\subset B$.

\subsection*{Correction des exercices de la section 6}
\addcontentsline{toc}{subsection}{Correction des exercices de la section 6}

\paragraph{Exercice \ref{exo:exo-inj-surj}.}
\begin{enumerate}
  \item $f:\R\to\R$, $f(x)=x^2$ n'est pas injective car $f(1)=f(-1)$ mais $1\neq-1$. Elle n'est pas surjective car $-1$ n'a pas d'antécédent réel.
  \item $g:\R_+\to\R_+$, $g(x)=x^2$ est injective : si $x^2=y^2$ avec $x,y\geq0$, alors $x=y$. Elle est surjective : pour $a\in\R_+$, $\sqrt a\in\R_+$ et $g(\sqrt a)=a$. Donc $g$ est bijective.
  \item $h(x)=x^3+x$ est injective car si $x<y$, alors $h(y)-h(x)=(y-x)(x^2+xy+y^2+1)>0$. Elle est surjective car $h$ est continue, $h(x)\to-\infty$ quand $x\to-\infty$ et $h(x)\to+\infty$ quand $x\to+\infty$ ; par le théorème des valeurs intermédiaires, tout réel est atteint.
\end{enumerate}

\paragraph{Exercice \ref{exo:exo-images}.}
Pour tout $y\in F$,
\begin{align*}
y\in f(A\cup B)
&\Longleftrightarrow \exists x\in A\cup B,\ y=f(x)\\
&\Longleftrightarrow (\exists x\in A,\ y=f(x))\text{ ou }(\exists x\in B,\ y=f(x))\\
&\Longleftrightarrow y\in f(A)\cup f(B).
\end{align*}
Donc $f(A\cup B)=f(A)\cup f(B)$.

Si $y\in f(A\cap B)$, alors il existe $x\in A\cap B$ tel que $y=f(x)$. Donc $x\in A$ et $x\in B$, d'où $y\in f(A)$ et $y\in f(B)$. Ainsi $y\in f(A)\cap f(B)$. Donc $f(A\cap B)\subset f(A)\cap f(B)$.

L'égalité peut être fausse. Prenons $f:\R\to\R$, $f(x)=x^2$, $A=\{-1\}$ et $B=\{1\}$. Alors $A\cap B=\varnothing$, donc $f(A\cap B)=\varnothing$, mais $f(A)=f(B)=\{1\}$, donc $f(A)\cap f(B)=\{1\}$.

\subsection*{Correction des exercices de la section 7}
\addcontentsline{toc}{subsection}{Correction des exercices de la section 7}

\paragraph{Exercice \ref{exo:exo-corriger-preuve}.}
L'erreur est que $x^2=1$ n'équivaut pas à $x=1$, car $x=-1$ est aussi solution. La résolution correcte est :
\[
x^2=1 \Longleftrightarrow x^2-1=0 \Longleftrightarrow (x-1)(x+1)=0 \Longleftrightarrow x=1\text{ ou }x=-1.
\]

\paragraph{Exercice \ref{exo:exo-redaction-equiv}.}
Montrons les deux implications.

Si $n$ est pair, alors il existe $k\in\Z$ tel que $n=2k$. Donc $n^2=4k^2=2(2k^2)$, donc $n^2$ est pair.

Réciproquement, supposons que $n^2$ est pair. Montrons que $n$ est pair par contraposée. Si $n$ est impair, alors $n=2k+1$ pour un certain $k\in\Z$, donc
\[
n^2=4k^2+4k+1=2(2k^2+2k)+1,
\]
ce qui est impair. Ainsi, si $n^2$ est pair, $n$ est pair. Les deux implications donnent l'équivalence.

%===========================================================
\section{Grand devoir bilan final}

\begin{center}
\begin{tcolorbox}[width=0.95\textwidth,colframe=bleuprepa,colback=blue!3!white,title={Devoir bilan -- Logique, raisonnement, ensembles et applications}]
Durée conseillée : 2 heures. Les réponses doivent être rédigées. Les symboles logiques peuvent être utilisés, mais toute démonstration doit être accompagnée de phrases claires.
\end{tcolorbox}
\end{center}

\begin{exercicebox}{Logique et tables de vérité}{bilan1}
\begin{enumerate}
  \item Construire la table de vérité de $(P\Rightarrow Q)\Longleftrightarrow(\neg Q\Rightarrow\neg P)$.
  \item Déterminer si $((P\Rightarrow Q)\land(Q\Rightarrow R))\Rightarrow(P\Rightarrow R)$ est une tautologie.
  \item Nier proprement : « si tous les élèves travaillent, alors il existe un élève qui réussit ».
\end{enumerate}
\end{exercicebox}

\begin{exercicebox}{Quantificateurs}{bilan2}
On considère la proposition
\[
\forall \varepsilon>0,\exists N\in\N,\forall n\in\N,\ n\geq N\Rightarrow |u_n-\ell|<\varepsilon.
\]
\begin{enumerate}
  \item Traduire cette proposition en français.
  \item La nier sans laisser de négation devant une proposition composée.
  \item Expliquer pourquoi l'ordre des quantificateurs $\forall\varepsilon\exists N$ est essentiel.
\end{enumerate}
\end{exercicebox}

\begin{exercicebox}{Raisonnements}{bilan3}
\begin{enumerate}
  \item Montrer que, pour tout $n\in\Z$, si $n^2$ est divisible par $3$, alors $n$ est divisible par $3$.
  \item Montrer par récurrence que pour tout $n\geq1$, $2^n\geq n+1$.
  \item Montrer qu'il n'existe pas de rationnel $r$ tel que $r^2=3$.
\end{enumerate}
\end{exercicebox}

\begin{exercicebox}{Ensembles et applications}{bilan4}
Soit $f:E\to F$ une application et soient $A,B\subset E$.
\begin{enumerate}
  \item Montrer que $f(A\cup B)=f(A)\cup f(B)$.
  \item Montrer que $f(A\cap B)\subset f(A)\cap f(B)$.
  \item Montrer que si $f$ est injective, alors $f(A\cap B)=f(A)\cap f(B)$.
\end{enumerate}
\end{exercicebox}

\section{Correction détaillée du devoir bilan final}

\paragraph{Correction du bilan 1.}
La première table montre que l'implication et sa contraposée ont toujours la même valeur de vérité. C'est donc une tautologie. Pour la deuxième, on reconnaît la transitivité de l'implication : si $P$ implique $Q$ et si $Q$ implique $R$, alors $P$ implique $R$. Une table de vérité complète confirme que la proposition est toujours vraie.

La phrase « si tous les élèves travaillent, alors il existe un élève qui réussit » a la forme $P\Rightarrow Q$, où $P$ signifie « tous les élèves travaillent » et $Q$ signifie « il existe un élève qui réussit ». Sa négation est $P\land\neg Q$ : « tous les élèves travaillent et aucun élève ne réussit ».

\paragraph{Correction du bilan 2.}
La proposition signifie : la suite $(u_n)$ converge vers $\ell$. Plus explicitement : pour toute précision $\varepsilon>0$, il existe un rang $N$ à partir duquel tous les termes de la suite sont à distance strictement inférieure à $\varepsilon$ de $\ell$.

Sa négation est :
\[
\exists \varepsilon>0,\forall N\in\N,\exists n\in\N,\ n\geq N\text{ et } |u_n-\ell|\geq\varepsilon.
\]
L'ordre des quantificateurs est essentiel : $N$ dépend de $\varepsilon$. Plus on demande une précision fine, plus il peut être nécessaire d'aller loin dans la suite.

\paragraph{Correction du bilan 3.}
\begin{enumerate}
  \item On raisonne par contraposée. Si $n$ n'est pas divisible par $3$, alors $n\equiv1$ ou $n\equiv2\pmod 3$. Dans les deux cas, $n^2\equiv1\pmod3$. Donc $n^2$ n'est pas divisible par $3$. La contraposée est démontrée.
  \item Pour $n=1$, $2^1=2\geq2$. Supposons $2^n\geq n+1$. Alors $2^{n+1}=2\cdot2^n\geq2(n+1)=2n+2\geq n+2$ pour $n\geq0$. Donc la propriété est héréditaire. Par récurrence, elle est vraie pour tout $n\geq1$.
  \item Supposons par l'absurde qu'il existe $r=p/q\in\Q$, avec $p,q$ premiers entre eux et $q\neq0$, tel que $r^2=3$. Alors $p^2=3q^2$, donc $3$ divise $p^2$, donc $3$ divise $p$. Écrivons $p=3k$. Alors $9k^2=3q^2$, donc $q^2=3k^2$, donc $3$ divise $q$. Contradiction avec le fait que $p$ et $q$ sont premiers entre eux. Donc aucun rationnel n'a pour carré $3$.
\end{enumerate}

\paragraph{Correction du bilan 4.}
Les deux premières questions ont été démontrées dans la correction de l'exercice sur les images. Pour la troisième, supposons $f$ injective. On sait déjà que $f(A\cap B)\subset f(A)\cap f(B)$. Montrons l'autre inclusion. Soit $y\in f(A)\cap f(B)$. Alors il existe $a\in A$ et $b\in B$ tels que $y=f(a)=f(b)$. Comme $f$ est injective, $a=b$. Cet élément commun appartient donc à $A\cap B$, et $y=f(a)\in f(A\cap B)$. Ainsi $f(A)\cap f(B)\subset f(A\cap B)$, d'où l'égalité.

%===========================================================
\section{Fiche de synthèse finale}

\begin{retenirbox}
\textbf{1. Une proposition} est une phrase vraie ou fausse. Un prédicat devient une proposition après fixation ou quantification des variables.
\end{retenirbox}

\begin{retenirbox}
\textbf{2. Connecteurs.} La négation inverse la vérité ; $P\land Q$ exige les deux ; $P\lor Q$ signifie au moins l'un des deux ; $P\Rightarrow Q$ est fausse seulement si $P$ est vraie et $Q$ fausse ; $P\Longleftrightarrow Q$ signifie deux implications.
\end{retenirbox}

\begin{retenirbox}
\textbf{3. Négations fondamentales.}
\[
\neg(P\land Q)\equiv\neg P\lor\neg Q,
\qquad
\neg(P\lor Q)\equiv\neg P\land\neg Q,
\]
\[
\neg(P\Rightarrow Q)\equiv P\land\neg Q.
\]
\[
\neg(\forall x, P(x))\equiv\exists x,\neg P(x),
\qquad
\neg(\exists x, P(x))\equiv\forall x,\neg P(x).
\]
\end{retenirbox}

\begin{retenirbox}
\textbf{4. Méthodes de preuve.} Directe, contraposée, absurde, disjonction de cas, double implication, analyse-synthèse, récurrence, contre-exemple.
\end{retenirbox}

\begin{retenirbox}
\textbf{5. Ensembles.} Une inclusion $A\subset B$ se démontre en partant de $x\in A$ et en prouvant $x\in B$. Une égalité d'ensembles se démontre souvent par double inclusion.
\end{retenirbox}

\begin{retenirbox}
\textbf{6. Applications.} Injectivité : deux antécédents ayant même image sont égaux. Surjectivité : tout élément de l'ensemble d'arrivée possède un antécédent. Bijectivité : les deux à la fois.
\end{retenirbox}

\begin{retenirbox}
\textbf{7. Rédaction.} Toujours annoncer ce qui est fixé, supposé, construit et démontré. Une figure suggère, mais seule une preuve justifie.
\end{retenirbox}

\section*{Références pédagogiques consultées pour le cadrage}
\addcontentsline{toc}{section}{Références pédagogiques consultées pour le cadrage}
\begin{itemize}
  \item Bulletin officiel spécial n\textsuperscript{o}1 du 11 février 2021 : programmes des CPGE scientifiques.
  \item Programmes de mathématiques MPSI, PCSI, MP2I et documents de cadrage officiels.
  \item Plans de cours publics de classes préparatoires portant sur logique, ensembles, applications et raisonnement mathématique.
\end{itemize}

\end{document}
